\chapter{Liver Transplant}

\section*{Introduction \& Clinical Indications}
Liver transplantation is a highly complex and life-saving surgical procedure that involves the replacement of a severely diseased or non-functional liver with a healthy liver from either a deceased organ donor or a segment of a liver from a living donor. This major abdominal operation is the definitive treatment for patients suffering from end-stage liver disease, acute liver failure, or specific liver cancers that are unresponsive to other therapeutic modalities. The primary goal of the surgery is to restore vital hepatic functions, including metabolism, detoxification, and protein synthesis, thereby offering a chance for extended survival and a significantly improved quality of life. It represents a transformative intervention for irreversible liver failure, fundamentally altering the prognosis for many patients who would otherwise face imminent mortality.

\begin{itemize}
  \item Hepatic Transplant
  \item Orthotopic Liver Transplant (OLT)
  \item Deceased Donor Liver Transplant (DDLT)
  \item Living Donor Liver Transplant (LDLT)
  \item Auxiliary Liver Transplant (rare, partial native liver preservation)
\end{itemize}

The fundamental surgical principle of orthotopic liver transplantation (OLT) is the complete excision of the recipient's diseased native liver and its replacement with a healthy donor liver, which is then meticulously revascularized and re-established with biliary drainage in the same anatomical position. This orthotopic placement ensures physiological blood flow and bile drainage, mimicking the natural architecture. The primary technical goals include achieving secure and patent vascular anastomoses for the hepatic artery, portal vein, and hepatic veins (or vena cava), and establishing effective biliary continuity to prevent complications.

The rationale for this procedure is to definitively treat irreversible liver failure, which can arise from various etiologies such as cirrhosis, acute fulminant hepatitis, or unresectable primary liver cancers. By replacing the non-functional organ, the transplant restores the liver's critical metabolic, synthetic, and detoxification functions, which are essential for life. For living donor liver transplantation (LDLT), a segment of the donor's liver (typically the right lobe for adults or left lateral segment for children) is used. The remarkable regenerative capacity of the liver allows both the donor's remaining liver and the transplanted segment in the recipient to grow back to near-normal size and function within weeks to months, providing a viable alternative to deceased donor organs and expanding the donor pool.

The journey of liver transplantation began in the early 1960s, marked by pioneering efforts and significant challenges. Dr. Thomas Starzl performed the first human liver transplant in Denver, Colorado, in 1963. While a technical success, the patient succumbed to rejection shortly after due to the nascent state of immunosuppression. Despite these initial setbacks, Starzl and his team persevered, achieving the first successful long-term liver transplant in 1967, a monumental milestone that demonstrated the procedure's potential.

The 1970s saw continued refinement of surgical techniques, but widespread success remained elusive due to high rates of graft rejection and severe surgical complications. A pivotal breakthrough occurred in the early 1980s with the introduction of cyclosporine, a potent calcineurin inhibitor. This immunosuppressive drug dramatically improved graft survival rates by effectively controlling rejection, revolutionizing transplant medicine. Subsequent developments, including the introduction of other immunosuppressants (e.g., tacrolimus, mycophenolate mofetil), refined organ preservation solutions (e.g., University of Wisconsin solution), and advancements in critical care management, firmly established liver transplantation as a viable and life-saving treatment by the mid-1980s, leading to its broader acceptance and application worldwide. Further evolution includes the development of living donor transplantation and split liver transplantation, expanding the donor pool and reducing waiting list mortality.

\begin{itemize}
  \item End-stage liver disease (ESLD) due to chronic viral hepatitis (Hepatitis B or C) leading to decompensated cirrhosis, characterized by ascites, encephalopathy, or variceal bleeding.
  \item Decompensated alcoholic liver disease, typically after a mandatory period of sustained abstinence and rehabilitation, demonstrating commitment to sobriety.
  \item Primary biliary cholangitis (PBC) or primary sclerosing cholangitis (PSC) when medical therapies fail and progressive cholestasis leads to intractable pruritus, recurrent cholangitis, or liver failure.
  \item Acute liver failure (ALF) from various etiologies, such as acetaminophen overdose, idiosyncratic drug reactions, or acute viral hepatitis, where rapid deterioration necessitates urgent intervention, often guided by King's College Criteria.
  \item Hepatocellular carcinoma (HCC) that meets specific criteria, such as the Milan criteria (single lesion \textless 5 cm or up to three lesions, none \textgreater 3 cm, without vascular invasion or extrahepatic spread), offering the best chance for long-term survival.
  \item Metabolic liver diseases, including alpha-1 antitrypsin deficiency, Wilson's disease, or hemochromatosis, which cause progressive and irreversible liver damage despite medical management.
  \item Certain congenital liver diseases in children, most notably biliary atresia, which leads to severe cholestasis and liver failure if not corrected early.
  \item Intractable complications of portal hypertension, such as recurrent variceal hemorrhage or hepatorenal syndrome, unresponsive to medical management and significantly impacting quality of life.
\end{itemize}

Liver transplantation is almost universally considered the primary and definitive treatment for end-stage liver disease and acute liver failure when other medical or surgical interventions have failed or are deemed insufficient to prevent progressive morbidity and mortality. For conditions like decompensated cirrhosis, it is the only curative option available that addresses the underlying organ failure. In cases of hepatocellular carcinoma, it serves as a primary oncologic treatment within strict selection criteria, offering superior long-term survival compared to other therapies for eligible patients, particularly those meeting Milan criteria. While some patients may undergo palliative treatments or bridge therapies, such as transjugular intrahepatic portosystemic shunt (TIPS) for refractory ascites or locoregional therapies (e.g., transarterial chemoembolization, radiofrequency ablation) for HCC, these are generally considered temporary measures to stabilize the patient or control disease progression while awaiting a suitable donor organ. Thus, liver transplantation is rarely a secondary or salvage procedure in the traditional sense, but rather the ultimate therapeutic goal for patients with irreversible liver pathology, offering the best chance for long-term survival and improved quality of life.

\section*{Relevant Surgical Anatomy}
The liver is the largest solid organ in the body, situated in the right upper quadrant of the abdomen, beneath the diaphragm, and protected by the lower ribs. It is broadly divided into right and left lobes, further subdivided into eight functional segments (Couinaud classification), each with its own vascular inflow, outflow, and biliary drainage. Key anatomical structures critical for liver transplantation include:
\begin{itemize}
  \item \textbf{Vascular Supply:} The liver receives a dual blood supply. The hepatic artery, typically arising from the celiac axis, provides approximately 25\% of the blood flow, rich in oxygen. The portal vein, formed by the confluence of the superior mesenteric and splenic veins, delivers about 75\% of the blood flow, rich in nutrients absorbed from the gastrointestinal tract. Venous drainage is via the hepatic veins (right, middle, and left), which drain directly into the inferior vena cava (IVC).
  \item \textbf{Biliary System:} Bile produced by hepatocytes drains into bile canaliculi, which coalesce into progressively larger ducts, forming the right and left hepatic ducts. These join to form the common hepatic duct, which then unites with the cystic duct (from the gallbladder) to form the common bile duct, emptying into the duodenum at the ampulla of Vater.
  \item \textbf{Nerve Relations:} The liver receives sympathetic innervation from the celiac plexus and parasympathetic innervation from the vagus nerve. These nerves primarily regulate blood flow and bile secretion.
  \item \textbf{Lymphatic Drainage:} Lymphatics generally follow the portal triads and drain into nodes around the porta hepatis, eventually reaching the celiac lymph nodes.
  \item \textbf{Surgical Landmarks:} The porta hepatis, containing the hepatic artery, portal vein, and common bile duct (forming the portal triad), is a crucial landmark for dissection and anastomosis. The retrohepatic IVC is intimately associated with the posterior surface of the liver, requiring careful dissection during hepatectomy.
\end{itemize}
Understanding these intricate relationships is paramount for safe and effective liver resection and transplantation, particularly during the complex vascular and biliary anastomoses.

\section*{Preoperative Workup \& Preparation}
Extensive preoperative evaluation is critical for both the recipient and, in living donor cases, the donor, to ensure they can withstand the physiological stress of surgery and the demands of lifelong immunosuppression. This multidisciplinary assessment aims to optimize the patient's condition and identify any contraindications.

Key preparatory steps include:
\begin{itemize}
  \item \textbf{Medical Evaluation:} Comprehensive assessment of cardiac function (ECG, echocardiogram, stress test, coronary angiography if indicated), pulmonary function (PFTs, chest X-ray), renal function, and nutritional status.
  \item \textbf{Infection Screening:} Thorough screening for active infections, including viral hepatitis (HBV, HCV, CMV, EBV), HIV, tuberculosis, and fungal infections. Any active infections must be treated or controlled prior to transplant.
  \item \textbf{Imaging Studies:} Abdominal CT or MRI to assess liver pathology, vascular anatomy, and rule out extrahepatic malignancy. Doppler ultrasound is crucial to evaluate portal vein patency and flow.
  \item \textbf{Psychosocial Evaluation:} Assessment of mental health, social support systems, and the patient's understanding and adherence to complex medical regimens, which are crucial for long-term success.
  \item \textbf{Immunological Workup:} Blood typing, tissue typing (HLA matching, though less critical than for kidney transplants), and cross-matching to minimize the risk of hyperacute rejection.
  \textbf{Listing and Allocation:} Once deemed eligible, the patient is listed on the national organ transplant registry (e.g., UNOS in the US) and allocated a MELD (Model for End-Stage Liver Disease) score or PELD (Pediatric End-Stage Liver Disease) score to prioritize organ allocation based on medical urgency.
  \item \textbf{Pre-operative Instructions:} Patients are typically NPO (nil per os) for several hours before surgery. Medications may be adjusted, particularly anticoagulants. Prophylactic broad-spectrum antibiotics are administered intravenously just prior to incision. Informed consent is obtained, detailing the extensive risks and benefits of the procedure.
\end{itemize}

\textbf{Situations requiring treatment, optimization, or center-specific review:}
\begin{itemize}
  \item Active, uncontrolled extrahepatic malignancy (outside the liver), as this indicates systemic disease and poor prognosis.
  \item Active, uncontrolled systemic infection or sepsis, which would lead to overwhelming infection post-transplant.
  \item Severe, irreversible cardiopulmonary disease that would preclude survival of surgery or recovery, such as severe pulmonary hypertension or advanced coronary artery disease unresponsive to treatment.
  \item Active, untreated substance-use disorder or inability to participate in an adherence plan; abstinence and rehabilitation requirements vary by center and clinical context.
  \item Documented non-adherence to medical therapy or follow-up, indicating poor likelihood of adherence to the complex post-transplant immunosuppressive regimen.
  \item Irreversible brain damage or other severe neurological deficits that would prevent meaningful recovery or participation in care.
  \item Inadequate ability to adhere to treatment or follow-up despite available support and interventions; social support should be assessed and strengthened rather than treated as a universal absolute exclusion.
\end{itemize}

\textbf{Relative Contraindications \& Precautions:}
\begin{itemize}
  \item Advanced age (generally \textgreater 70 years), though physiological age and overall health are more important than chronological age, assessed on a case-by-case basis.
  \item Significant obesity (BMI \textgreater 35-40 kg/m\textasciicircum2) due to increased surgical risk, technical challenges, and higher rates of post-operative complications.
  \item Extensive portal vein thrombosis that is non-reconstructible, making portal inflow difficult to establish.
  \item Prior extensive upper abdominal surgery, which can increase the complexity and duration of the hepatectomy due to adhesions.
  \item HIV infection, though with effective antiretroviral therapy and controlled viral load, it is becoming less of a contraindication in selected patients.
  \item Active psychiatric illness that is not well-controlled, which may impair adherence to the post-transplant regimen.
  \item Hepatocellular carcinoma exceeding Milan criteria, though some centers use expanded criteria for selected patients with careful risk-benefit assessment.
\end{itemize}

\section*{Surgical Technique \& Procedure}
Liver transplantation is a highly complex surgical procedure requiring a multidisciplinary team and meticulous execution.
\begin{enumerate}
  \item \textbf{Anesthesia and Positioning:} The patient is placed under general anesthesia with extensive hemodynamic monitoring, including central venous and arterial lines. They are positioned supine, and a large incision, typically a bilateral subcostal incision with a midline extension (Mercedes incision) or a modified J-incision, is made to provide optimal exposure of the abdominal cavity and liver.
  \item \textbf{Hepatectomy Phase (Recipient):} The diseased native liver is meticulously dissected free from its surrounding attachments. This involves ligating and dividing the hepatic artery, portal vein, and bile duct. The retrohepatic vena cava is often preserved using the "piggyback" technique, where the recipient's hepatic veins are anastomosed directly to the donor's hepatic veins, leaving the recipient's IVC intact. In other cases, a full caval replacement may be necessary, requiring temporary venovenous bypass to maintain hemodynamic stability. Careful hemostasis is paramount due to the highly vascular nature of the liver and the presence of portal hypertension.
  \item \textbf{Anhepatic Phase (if full caval clamping):} This is the critical period after the native liver is removed and before the donor liver is fully revascularized. It is associated with significant physiological challenges, including acidosis, hypothermia, and coagulopathy, requiring careful anesthetic and surgical management.
  \item \textbf{Implantation and Reperfusion Phase (Donor Liver):} The donor liver, preserved in cold storage solution, is brought to the field. The major vascular anastomoses are performed first:
    \begin{itemize}
      \item Suprahepatic vena cava anastomosis.
      \item Infrahepatic vena cava anastomosis (if full caval replacement) or direct hepatic vein anastomoses to the recipient vena cava (piggyback technique).
      \item Portal vein anastomosis.
      \item Hepatic artery anastomosis, often using microsurgical techniques due to its small caliber.
    \end{itemize}
    Once vascular anastomoses are complete, the clamps are released, allowing blood flow to the new liver (reperfusion). This can lead to "reperfusion syndrome," characterized by transient hypotension and arrhythmias, requiring careful management.
  \item \textbf{Biliary Reconstruction:} Finally, biliary continuity is restored. The most common method is a duct-to-duct anastomosis between the donor and recipient common bile ducts. If the recipient's bile duct is diseased, scarred, or unsuitable, a Roux-en-Y choledochojejunostomy may be performed, connecting the donor bile duct to a loop of recipient jejunum.
  \item \textbf{Closure:} Once all anastomoses are complete and hemostasis is secured, surgical drains are often placed to monitor for fluid collections or bleeding. The abdomen is then closed in layers.
\end{enumerate}
The entire procedure can take 6-12 hours, depending on the complexity and the presence of adhesions from previous surgeries or severe portal hypertension.

\section*{Complications \& Risks}
Liver transplantation carries significant risks, both immediate and long-term, due to the complexity of the surgery and the necessity for lifelong immunosuppression.

\begin{itemize}
  \item \textbf{General Surgical Complications:}
    \begin{itemize}
      \item \textbf{Bleeding:} Significant intraoperative and postoperative blood loss requiring transfusions is common due to portal hypertension, coagulopathy, and extensive dissection.
      \item \textbf{Infection:} High risk of bacterial, viral (e.g., CMV, EBV), and fungal infections, both surgical site and systemic, exacerbated by immunosuppression.
      \item \textbf{Anesthesia Complications:} Reactions to medications, respiratory or cardiac arrest, stroke, or myocardial infarction.
      \item \textbf{Wound Complications:} Incisional hernia, wound dehiscence, or infection.
      \item \textbf{Thromboembolic Events:} Deep vein thrombosis (DVT) and pulmonary embolism (PE).
    \end{itemize}
  \item \textbf{Procedure-Specific Risks:}
    \begin{itemize}
      \item \textbf{Primary Graft Non-Function (PNF):} The transplanted liver fails to function adequately immediately after surgery, often requiring urgent re-transplantation.
      \item \textbf{Vascular Complications:}
        \begin{itemize}
          \item \textit{Hepatic Artery Thrombosis (HAT):} A critical complication (incidence 2-9\%) leading to graft ischemia, biliary necrosis, and potential failure, often requiring urgent re-transplantation.
          \item \textit{Portal Vein Thrombosis:} Can impair graft perfusion and function, potentially leading to portal hypertension.
          \item \textit{Vena Cava Stenosis/Thrombosis:} Can lead to lower extremity edema, renal dysfunction, or Budd-Chiari-like syndrome.
        \end{itemize}
      \item \textbf{Biliary Complications:}
        \begin{itemize}
          \item \textit{Bile Leak:} From the anastomotic site, leading to peritonitis, abscess formation, or external fistula.
          \item \textit{Biliary Stricture:} Narrowing of the bile duct anastomosis or intrahepatic ducts, causing cholestasis, cholangitis, and potentially graft dysfunction, often requiring endoscopic or percutaneous intervention.
        \end{itemize}
      \item \textbf{Rejection:} The recipient's immune system attacks the donor liver. Can be acute (cellular or antibody-mediated) or chronic, requiring adjustment of immunosuppression and potentially leading to graft loss.
      \item \textbf{Neurological Complications:} Post-transplant encephalopathy, seizures, stroke, or central pontine myelinolysis.
      \item \textbf{Renal Dysfunction:} Can be exacerbated by calcineurin inhibitors (immunosuppressants), potentially leading to chronic kidney disease.
      \item \textbf{Recurrence of Original Disease:} Particularly for viral hepatitis (HCV, HBV), autoimmune hepatitis, or hepatocellular carcinoma, requiring close monitoring and specific antiviral or oncologic therapies.
      \item \textbf{Immunosuppression-Related Complications:} Increased risk of opportunistic infections, new onset diabetes mellitus, hypertension, hyperlipidemia, osteoporosis, and certain malignancies (e.g., post-transplant lymphoproliferative disorder (PTLD), skin cancers, kidney cancer).
    \end{itemize}
\end{itemize}

\section*{Postoperative Care \& Recovery}
Immediately following liver transplantation, the patient is transferred to the intensive care unit (ICU) for intensive monitoring. This critical period involves continuous assessment of vital signs, hemodynamic stability, urine output, and neurological status. Liver function tests (LFTs) are monitored frequently (every 4-6 hours initially) to assess graft function, and Doppler ultrasound scans are performed regularly to ensure patency of the hepatic artery and portal vein. Pain management is initiated with intravenous analgesics, and immunosuppressive medications are started to prevent rejection. The first 24-48 hours are crucial for stabilizing the patient, managing fluid and electrolyte shifts, and detecting early complications like primary graft non-function or vascular thrombosis.

Over the first few days to weeks, if the graft is functioning well and the patient is stable, they will be transferred from the ICU to a specialized transplant ward. During this time, the focus shifts to optimizing immunosuppression, managing fluid and electrolyte balance, preventing infections through prophylactic antibiotics and antivirals, and initiating early mobilization. Diet is gradually advanced from clear liquids to a regular, healthy diet as tolerated. Patients typically remain hospitalized for 2-4 weeks, depending on their recovery trajectory and any complications. Before discharge, extensive education regarding medication adherence, signs of rejection or infection, dietary restrictions, and activity limitations is provided. Long-term recovery involves lifelong immunosuppression and regular follow-up with the transplant team to monitor for rejection, infection, recurrence of original disease, and side effects of immunosuppressive medications, allowing most patients to return to a productive life within 3-6 months.

During the hepatectomy phase, the surgeon meticulously documents the gross appearance of the native liver, noting the extent of cirrhosis, fibrosis, nodularity, presence of ascites, varices, and any tumors (e.g., hepatocellular carcinoma). The size and consistency of the liver, as well as the presence of adhesions from previous surgeries or inflammation, are also recorded. Any unexpected findings, such as extrahepatic malignancy or significant vascular anomalies, are noted and may influence subsequent management.

The resected native liver is sent to pathology for detailed macroscopic and microscopic examination. The pathology report confirms the underlying liver disease (e.g., type and severity of cirrhosis, presence of viral hepatitis, alcoholic liver disease, autoimmune hepatitis, or metabolic disorders). For patients with HCC, the report details the tumor size, number, differentiation, presence of vascular invasion, and margin status, which are crucial for assessing the risk of recurrence. The donor liver is also subject to pathological assessment, primarily to confirm its health and rule out any unexpected abnormalities that could compromise graft function. This comprehensive pathological analysis provides vital information for understanding the disease progression, confirming the indication for transplant, and guiding long-term postoperative management and surveillance.

A normal, successful postoperative course after liver transplantation is characterized by a rapid and sustained improvement in the patient's clinical condition and liver function. Within days to weeks, pre-transplant symptoms such as jaundice, ascites, and encephalopathy should progressively resolve. Laboratory tests should show a rapid decline in serum bilirubin levels, normalization of coagulation parameters (INR, PT), and improvement in synthetic function (e.g., rising albumin). While elevated liver enzymes (AST, ALT) are common initially due to reperfusion injury, they should trend downwards. The patient should experience progressive improvement in energy levels, appetite, and overall well-being, allowing for early mobilization and rehabilitation. Pain should be well-controlled and gradually diminish, transitioning from intravenous to oral analgesics. Imaging studies, particularly Doppler ultrasound, should consistently confirm patent hepatic artery and portal vein flow, without evidence of thrombosis or significant biliary dilatation. Ultimately, an expected normal course signifies a stable, functioning liver graft with minimal complications, enabling the patient to resume a productive and active life under lifelong medical surveillance.

Lifelong follow-up is essential after liver transplantation to ensure graft function, manage immunosuppression, and detect complications early.
\begin{itemize}
  \item \textbf{Regular Clinic Visits:} Frequent visits in the first year (weekly to monthly) gradually decreasing to annual visits thereafter, with the transplant team (hepatologist, surgeon, coordinator).
  \item \textbf{Blood Tests:}
    \begin{itemize}
      \item \textit{Immunosuppressant Drug Levels:} Monitored frequently to maintain therapeutic ranges and minimize toxicity (e.g., tacrolimus, cyclosporine).
      \item \textit{Liver Function Tests (LFTs):} To assess graft health and detect early signs of dysfunction or rejection.
      \item \textit{Kidney Function Tests:} To monitor for nephrotoxicity from immunosuppressants.
      \item \textit{Complete Blood Count (CBC):} To check for anemia, leukopenia, or thrombocytopenia.
      \item \textit{Infection Markers:} Regular screening for viral infections (CMV, EBV, HBV, HCV) and other opportunistic pathogens.
      \item \textit{Metabolic Panel:} To monitor for new onset diabetes, hyperlipidemia, and electrolyte imbalances.
    \end{itemize}
  \item \textbf{Imaging Studies:}
    \begin{itemize}
      \item \textit{Doppler Ultrasound:} Periodically to assess hepatic artery and portal vein patency and detect biliary dilatation.
      \item \textit{CT or MRI:} As needed to investigate specific concerns or screen for recurrence of malignancy.
    \end{itemize}
  \item \textbf{Biopsies:} Liver biopsies may be performed if there are unexplained elevations in LFTs to diagnose rejection or other graft pathology.
  \item \textbf{Endoscopic Procedures:} ERCP (Endoscopic Retrograde Cholangiopancreatography) or PTC (Percutaneous Transhepatic Cholangiography) may be required to manage biliary strictures or leaks.
  \item \textbf{Cancer Screening:} Increased surveillance for post-transplant malignancies, including skin cancer, PTLD, and recurrence of HCC.
  \item \textbf{Vaccinations:} Regular vaccinations are important, though live vaccines are generally contraindicated due to immunosuppression.
  \item \textbf{Lifestyle Management:} Adherence to a healthy diet, regular exercise, and avoidance of alcohol and smoking are strongly encouraged.
\end{itemize}
Patients are educated on warning signs such as fever, jaundice, dark urine, light stools, severe abdominal pain, or sudden fatigue, which warrant immediate medical attention.

\section*{Outcomes \& Prognosis}
Liver transplantation is a relatively common major organ transplant procedure globally, with over 8,000 transplants performed annually in the United States alone and increasing numbers worldwide. The incidence of end-stage liver disease, the primary indication for transplant, varies geographically but is rising in many regions due to increasing rates of non-alcoholic fatty liver disease (NAFLD) and its more severe form, non-alcoholic steatohepatitis (NASH), which is projected to become the leading indication. While historically chronic viral hepatitis (Hepatitis C) was a leading cause, successful antiviral therapies have shifted the landscape, with alcoholic liver disease and NAFLD/NASH now representing increasingly frequent indications. Liver transplantation is performed across all age groups, from infants with congenital conditions like biliary atresia to elderly patients with advanced cirrhosis, though the majority of recipients are adults between 40 and 70 years of age. There is a slight male predominance among recipients, largely reflecting the higher incidence of alcoholic liver disease and hepatocellular carcinoma in men. Despite advancements, the mortality rate for patients on the transplant waiting list remains significant, underscoring the critical shortage of donor organs. Post-transplant morbidity is primarily related to immunosuppression-related complications and graft dysfunction, with infection and rejection being major contributors.

Outcomes are substantially better than without transplantation for appropriately selected patients, but they vary by indication, urgency, donor and graft factors, comorbidity, center, and post-transplant complications. Prognosis is influenced by graft function, rejection, infection, vascular or biliary complications, recurrence of the original disease, renal function, and adherence to immunosuppression. Lifelong follow-up and immunosuppression are usually required, although regimens are individualized and some patients may later be considered for carefully supervised minimization.

\section*{References}
\begin{enumerate}
  \item Schwartz's Principles of Surgery. 11th ed. New York, NY: McGraw-Hill Education; 2019. Chapter 32: Liver Transplantation.
  \item Sabiston Textbook of Surgery: The Biological Basis of Modern Surgical Practice. 21st ed. Philadelphia, PA: Elsevier; 2021. Chapter 56: Liver Transplantation.
  \item European Association for the Study of the Liver (EASL). Clinical Practice Guidelines: Liver Transplantation. J Hepatol. 2016;64(2):433-485.
  \item MedlinePlus. Liver transplant. U.S. National Library of Medicine; 2024. Available from: \url{https://medlineplus.gov/livertransplant.html}
\end{enumerate}

\clearpage
